\documentclass[11pt]{article}
\usepackage[margin=24mm]{geometry}\usepackage[T1]{fontenc}
\usepackage{amsmath,amssymb,amsthm,mathtools,mathrsfs,longtable,booktabs,array,graphicx,xurl}
\usepackage[hidelinks,hypertexnames=false]{hyperref}
\allowdisplaybreaks\setlength{\emergencystretch}{3em}
\providecommand{\tightlist}{\setlength{\itemsep}{0pt}\setlength{\parskip}{0pt}}
\DeclareUnicodeCharacter{03C0}{\ensuremath{\pi}}
\DeclareUnicodeCharacter{03C7}{\ensuremath{\chi}}
\DeclareUnicodeCharacter{03A6}{\ensuremath{\Phi}}
\DeclareUnicodeCharacter{0394}{\ensuremath{\Delta}}
\DeclareUnicodeCharacter{0393}{\ensuremath{\Gamma}}
\DeclareUnicodeCharacter{039B}{\ensuremath{\Lambda}}
\DeclareUnicodeCharacter{03C3}{\ensuremath{\sigma}}
\DeclareUnicodeCharacter{03B7}{\ensuremath{\eta}}
\DeclareUnicodeCharacter{03C4}{\ensuremath{\tau}}
\DeclareUnicodeCharacter{03B4}{\ensuremath{\delta}}
\DeclareUnicodeCharacter{03B3}{\ensuremath{\gamma}}
\DeclareUnicodeCharacter{03B1}{\ensuremath{\alpha}}
\DeclareUnicodeCharacter{03B2}{\ensuremath{\beta}}
\DeclareUnicodeCharacter{03C1}{\ensuremath{\rho}}
\DeclareUnicodeCharacter{03BA}{\ensuremath{\kappa}}
\DeclareUnicodeCharacter{03BB}{\ensuremath{\lambda}}
\DeclareUnicodeCharacter{03B8}{\ensuremath{\theta}}
\DeclareUnicodeCharacter{03B5}{\ensuremath{\varepsilon}}
\DeclareUnicodeCharacter{2208}{\ensuremath{\in}}
\DeclareUnicodeCharacter{2264}{\ensuremath{\le}}
\DeclareUnicodeCharacter{2265}{\ensuremath{\ge}}
\DeclareUnicodeCharacter{2260}{\ensuremath{\ne}}
\DeclareUnicodeCharacter{2192}{\ensuremath{\to}}
\DeclareUnicodeCharacter{2282}{\ensuremath{\subset}}
\DeclareUnicodeCharacter{2286}{\ensuremath{\subseteq}}
\DeclareUnicodeCharacter{2297}{\ensuremath{\otimes}}
\DeclareUnicodeCharacter{03A3}{\ensuremath{\Sigma}}
\DeclareUnicodeCharacter{221A}{\ensuremath{\surd}}
\DeclareUnicodeCharacter{221E}{\ensuremath{\infty}}
\DeclareUnicodeCharacter{00B2}{\ensuremath{{}^2}}
\DeclareUnicodeCharacter{00B3}{\ensuremath{{}^3}}
\DeclareUnicodeCharacter{2074}{\ensuremath{{}^4}}
\DeclareUnicodeCharacter{2080}{\ensuremath{{}_0}}
\DeclareUnicodeCharacter{2081}{\ensuremath{{}_1}}
\DeclareUnicodeCharacter{2082}{\ensuremath{{}_2}}
\DeclareUnicodeCharacter{2096}{\ensuremath{{}_k}}
\DeclareUnicodeCharacter{2212}{\ensuremath{-}}
\DeclareUnicodeCharacter{00D7}{\ensuremath{\times}}
\DeclareUnicodeCharacter{00B1}{\ensuremath{\pm}}
\DeclareUnicodeCharacter{2202}{\ensuremath{\partial}}
\DeclareUnicodeCharacter{2225}{\ensuremath{\Vert}}
\DeclareUnicodeCharacter{2113}{\ensuremath{\ell}}
\DeclareUnicodeCharacter{211D}{\ensuremath{\mathbb R}}
\DeclareUnicodeCharacter{2102}{\ensuremath{\mathbb C}}
\DeclareUnicodeCharacter{03BC}{\ensuremath{\mu}}
\DeclareUnicodeCharacter{03B6}{\ensuremath{\zeta}}
\DeclareUnicodeCharacter{03BE}{\ensuremath{\xi}}
\DeclareUnicodeCharacter{220F}{\ensuremath{\prod}}
\begin{document}
\section{Independent proof of the polynomial and tensor coercivity estimates}\label{independent-proof-of-the-polynomial-and-tensor-coercivity-estimates}

Date: 2026-09-20. Scope: equations P1--P11 and J1--J11 in the supplied \path{COMPLETE_PROOFS(2).md}, Sections 1--2. This file gives the complete independent derivation within that scope. Sections 3--5 have not been audited here.

\subsection{Acceptance and precise scope}\label{acceptance-and-precise-scope}

All numerical constants in P1--P11 are valid. J1--J11 are valid for the original nested sources, a nonempty fixed divisor, positive fixed jet weights, and a single section fixed at the four cutoffs. The formulas using the inverse of the scalar minimum metric require its attained domain, namely N \textgreater= q-1. The joint minimum itself exists for every N \textgreater= 0 because the section is onto.

One sentence requires a scope correction. The last paragraph of Section 2 extends all results to intermediate sources inside the same polynomial box. The common metric bounds extend to every such source containing the section. The nonnegative lower bound in J11 additionally requires that these sources increase with N. Without this nesting, the enclosure gives the valid bound \texttt{\textbar{}R\ log\ det\textbar{}\ \textless{}=\ 2\ r\_X\ W\_k}, but it does not prove a nonnegative return. The actual sources \texttt{X\_N\ +\ Y\_\{k,eta\}} are nested, so their stated J11 is valid.

For general multiplicities the exact tolerance cost is retained below. Nothing here removes \texttt{log(1/eta)} from that cost when eta is allowed to shrink arbitrarily. For simple factors, the metric has no higher jets and the cost does not depend on eta.

The displayed maxima over roots and the degree \texttt{n\_0=2d-1} in the supplied formulas presuppose a nonempty divisor. This is the domain used below. No convention for an empty root set is silently inserted.

\subsection{1. Original spaces, coordinates, and physical unit}\label{original-spaces-coordinates-and-physical-unit}

Let Z be a finite nonempty set of distinct nontrivial zeros of zeta, and let each m\_rho be its complete order in 2 xi. Put \[
h(s)=\prod_{\rho\in Z}(s-\rho)^{m_\rho},\quad
d=\sum_\rho m_\rho,\quad d_\rho=m_\rho-1,\quad
d_{\max}=\max_\rho d_\rho.
\] Fix a\_rho \textgreater{} 0. The physical line is s=1/2+iy. In all norm integrals, \[
\|P\|_{w_h}^2=\int_{\mathbb R}|P(1/2+iy)|^2
 \frac{|(2\xi/h)(1/2+iy)|^2}{2\pi}\,dy,
\] \[
\|P\|_\sigma^2=\int_{\mathbb R}|P(1/2+iy)|^2
 \frac{|\Gamma(1/4+iy/2)|^2}{2\pi}\,dy.
\] When a polynomial is displayed as a function of y, it means the exact substitution Q(y)=P(1/2+iy). No density is divided by its mass.

The coefficient quotient is A\_h=C{[}s{]}/(h), represented by polynomials of degree below d. The divided-jet map is \[
\mathsf H_h:A_h\longrightarrow
 \bigoplus_{\rho\in Z}\mathbb C[t_\rho]/(t_\rho^{m_\rho}),\qquad
[P]\longmapsto\left(\sum_{j=0}^{d_\rho}
 \frac{P^{(j)}(\rho)}{j!}t_\rho^j\right)_\rho.
\] This is an isomorphism: its kernel consists precisely of multiples of h, and its source and target both have dimension d. In these jets define \[
\mathsf D_{h,\eta}[P]=\sum_{\rho,j}
a_\rho(j!)^2\binom{d_\rho}{j}\eta^{2j}
\left|\frac{P^{(j)}(\rho)}{j!}\right|^2,
\quad 0<\eta\le1,
\qquad T_\eta=\mathsf H_h^*\mathsf D_{h,\eta}\mathsf H_h.
\] Here and below a positive matrix also denotes its quadratic form when the argument is written in brackets.

Put f\_h=2xi/h. It is entire and f\_h(rho) is nonzero at every selected root, because the complete order was removed. Multiplication by its truncated Taylor series defines the invertible jet map U\_h. Every coefficient is retained: \[
(U_h u)_{\rho,j}=\sum_{r=0}^j
 \frac{f_h^{(r)}(\rho)}{r!}\,u_{\rho,j-r}.
\] With the physical metric D\_phys=(U\_h\^{}\{-1\})\^{}*D\_\{h,1\}U\_h\^{}\{-1\}, \[
\|U_hu\|_{D_{\rm phys}}^2=u^*D_{h,1}u.
\] The tensor identity is obtained by taking tensor powers of this exact identity. Thus the later value constraint beta v and the physical constraint \(U_h^{\otimes k}\beta v\) have equal specified norms; the physical unit and its derivatives have not been discarded.

\subsection{2. Gamma polynomial basis, mass, and multiplication bound}\label{gamma-polynomial-basis-mass-and-multiplication-bound}

The Meixner--Pollaczek formulas used are DLMF 18.23.7, 18.19.8, and 18.19.9. Substitute lambda=1/4, phi=pi/2, x=y/2, and define \[
p_m(y)=m!P_m^{(1/4)}(y/2;\pi/2).
\] Then \[
\sum_{m\ge0}\frac{p_m(y)}{m!}z^m
 =(1+z^2)^{-1/4}\exp(y\arctan z),\qquad |z|<1.
\tag{C1}
\] The leading coefficient in y is one, as is also immediate by selecting the highest y power in C1. The exact change dy=2dx in DLMF\textquotesingle s orthogonality gives \[
\langle p_m,p_l\rangle_\sigma=\delta_{ml}\gamma_m,
\qquad
\gamma_m=\sqrt{2\pi}\,m!(1/2)_m.
\tag{C2}
\] Indeed the DLMF norm of P\_m is \texttt{2\ pi\ Gamma(m+1/2)/(sqrt(2)\ m!)}; multiplying by \texttt{(m!)\^{}2/pi} gives C2. In particular M\_sigma=gamma\_0=sqrt(2 pi), exactly as in P10.

Differentiate C1 in z. Its right side satisfies \[
(1+z^2)\partial_z C1=(y-z/2)C1.
\] Comparison of coefficients gives, with p\_\{-1\}=0, \[
yp_m=p_{m+1}+m(m-1/2)p_{m-1}.
\tag{C3}
\] Thus in the orthonormal basis e\_m=p\_m/sqrt(gamma\_m), the off-diagonal coefficient is alpha\_m=sqrt(m(m-1/2)). If deg Q \textless= n, the two shifted coefficient vectors in yQ have norms at most alpha\_\{n+1\}\textbar\textbar Q\textbar\textbar{} and alpha\_n\textbar\textbar Q\textbar\textbar. The triangle inequality yields \[
\|yQ\|_\sigma\le(\alpha_{n+1}+\alpha_n)\|Q\|_\sigma
 \le2(n+1)\|Q\|_\sigma.
\tag{C4}
\] For R=8(n+1), integration of \texttt{1\_\{\textbar{}y\textbar{}\textgreater{}R\}\ \textless{}=\ y\^{}2/R\^{}2} now gives \[
\int_{|y|>R}|Q(y)|^2\sigma(y)\,dy\le\frac1{16}\|Q\|_\sigma^2.
\tag{C5}
\] This verifies P5, including its degree and constant.

\subsection{3. Legendre localization and the excluded mass}\label{legendre-localization-and-the-excluded-mass}

DLMF 5.7.6 gives, for a\textgreater0 and t\textgreater=0, \[
\operatorname{Im}\psi(a+it)=\frac{t}{a^2+t^2}
 +\sum_{j\ge1}\frac{t}{(j+a)^2+t^2}.
\] The first summand is at most 1/(2a). The remaining summand is decreasing as a function of j\textgreater=0, so its sum for j\textgreater=1 is at most \[
\int_0^\infty\frac{t}{(x+a)^2+t^2}\,dx\le\pi/2.
\] At a=1/4 this proves \texttt{\textbar{}Im\ psi(a+it)\textbar{}\ \textless{}=\ 2+pi/2\ \textless{}4}. Differentiating the actual density, with its factor y/2, gives \[
\frac d{dy}\log\sigma(y)=-\operatorname{Im}\psi(1/4+iy/2).
\] Complex conjugation handles negative y. Therefore, for \textbar u-y\textbar\textless=1, \[
\sigma(y)\le e^4\sigma(u).
\tag{C6}
\]

Let L\_m be the Legendre polynomial on {[}-1,1{]}, with \texttt{integral\ L\_m\^{}2=2/(2m+1)}. Its value at zero is 0 for odd m; for m=2j it is \texttt{(-1)\^{}j\ binom(2j,j)/4\^{}j}, of modulus at most 1. Expand the complete complex polynomial Q(y+t) in the orthonormal functions \texttt{sqrt((2m+1)/2)L\_m(t)}. Cauchy--Schwarz at t=0 gives \[
|Q(y)|^2\le
 \left(\sum_{m=0}^n\frac{2m+1}{2}|L_m(0)|^2\right)
 \int_{y-1}^{y+1}|Q(u)|^2\,du
\le\frac{(n+1)^2}{2}\int_{y-1}^{y+1}|Q(u)|^2\,du.
\] Multiply by sigma(y), use C6 inside the integral, and extend that integral to the whole line. This proves P4 with exactly \[
B_n=e^4(n+1)^2/2.
\tag{C7}
\] Consequently every measurable excluded set E of length at most 1/(32B\_n) has \texttt{integral\_E\ \textbar{}Q\textbar{}\^{}2\ sigma\ \textless{}=\ \textbar{}\textbar{}Q\textbar{}\textbar{}\_sigma\^{}2/32}.

\subsection{4. The zeta bound and all disk constants}\label{the-zeta-bound-and-all-disk-constants}

Set g(s)=(s-1)zeta(s), which is entire. Taking n=1 in DLMF 25.2.10 and integrating its periodic B\_3 remainder using B\_4\textquotesingle=4B\_3 gives \[
\zeta(s)=\frac1{s-1}+\frac12+\frac{s}{12}
 -\frac{s(s+1)(s+2)}{720}
 -\frac{(s)_4}{24}\int_1^\infty\widetilde B_4(x)x^{-s-4}\,dx.
\tag{C8}
\] Initially this integration is valid on the common convergence domain; the resulting absolutely convergent integral extends the identity to Re s \textgreater{} -3. The sign and coefficient of its boundary term follow from \texttt{B\_4(1)=-1/30} and \texttt{{[}-B\_4(1)/4{]}=1/120} in the integration of B\_3. Since \texttt{B\_4(t)=t\^{}2(1-t)\^{}2-1/30} on {[}0,1{]}, its range lies in \texttt{{[}-1/30,7/240{]}}; hence \texttt{sup\ \textbar{}B\_4\textbar{}=1/30}.

If Re s \textgreater= -2, the absolute value of the integral in C8 is at most 1/30. Put t=\textbar s\textbar, T=t+4\textgreater=4. Multiplication by s-1 yields the explicit estimate \[
|g(s)|\le1+\frac{t+1}{2}+\frac{t(t+1)}{12}
 +\frac{t(t+1)^2(t+2)}{720}
 +\frac{t(t+1)^2(t+2)(t+3)}{720}
\le1+T/2+T^2/12+T^4/720+T^5/720<T^5.
\tag{C9}
\] For the last inequality, division by T\^{}5 bounds the right side by \texttt{1/1024+1/512+1/768+1/2880+1/720\ \textless{}1}. At the removable point s=1 the same inequality follows by continuity. This verifies P6.

The Euler product in Re s\textgreater1 gives \[
|\zeta(2+ij)^{-1}|\le\prod_p(1+p^{-2})
 \le\prod_p(1-p^{-2})^{-1}=\zeta(2)<2.
\] The last strict inequality also follows directly from the sum and the integral comparison \texttt{sum\_\{m\textgreater{}=2\}m\^{}\{-2\}\textless{}integral\_1\^{}infty\ x\^{}\{-2\}dx=1}. Since \textbar1+ij\textbar\textgreater=1, \texttt{\textbar{}g(2+ij)\textbar{}\textgreater{}1/2}.

Fix an integer j with \textbar j\textbar\textless=R, and put f\_j(z)=g(2+ij+z). On \textbar z\textbar\textless=4, Re(2+ij+z)\textgreater=-2 and \texttt{\textbar{}2+ij+z\textbar{}+4\ \textless{}=\ R+10}. Therefore C9 gives \[
|f_j(z)|\le(R+10)^5\le M_n:=4(R+14)^5.
\tag{C10}
\] Let L\_n=log(2M\_n), and let N\_j be the number of zeros a of f\_j in \textbar a\textbar\textless3, counted with multiplicity. Jensen\textquotesingle s formula at radii tending to 4 from below gives \[
N_j\log(4/3)\le\log M_n-\log|f_j(0)|<L_n.
\] No zero occurs at zero. Thus \[
N_j\le K_n=\left\lceil L_n/\log(4/3)\right\rceil.
\tag{C11}
\] This reasoning allows zeros on the circle of radius 4 by taking the limit through radii avoiding their moduli.

Form the finite product \[
B_j(z)=\prod_{|a|<3}\frac{3(z-a)}{9-\bar a z},
\qquad F_j(z)=f_j(z)/B_j(z),
\] with multiplicities. All apparent poles of F\_j at the listed zeros are removable, and there are no zeros of F\_j in \textbar z\textbar\textless3. On \textbar z\textbar=3 every factor has modulus one, so the maximum principle gives \texttt{\textbar{}F\_j\textbar{}\textless{}=M\_n} on the disk. Also \texttt{\textbar{}B\_j(0)\textbar{}\textless{}=1}, hence \texttt{\textbar{}F\_j(0)\textbar{}\textgreater{}1/2}.

The nonnegative harmonic function H=log(M\_n/\textbar F\_j\textbar) has H(0)\textless L\_n. To apply Harnack without any boundary regularity assumption, use every radius r with 2\textless r\textless3 and then let r increase to 3. At \textbar z\textbar\textless=2, \[
H(z)\le\lim_{r\uparrow3}\frac{r+2}{r-2}H(0)
 \le5L_n.
\] Thus \[
|F_j(z)|\ge M_ne^{-5L_n}=\frac1{32M_n^4},\qquad |z|\le2.
\tag{C12}
\] For each removed zero, \texttt{\textbar{}9-bar(a)z\textbar{}\textless{}=15} on \textbar z\textbar\textless=2. Hence \[
\left|\frac{3(z-a)}{9-\bar a z}\right|\ge|z-a|/5.
\tag{C13}
\]

The critical-line cell \texttt{j-1/2\ \textless{}=\ y\ \textless{}=\ j+1/2} has relative coordinate \texttt{z=-3/2+i(y-j)}, whose modulus is at most sqrt(5/2)\textless2. Let \[
m_n=2R+1,\qquad
\epsilon_n=(64B_nm_nK_n)^{-1},\qquad
c_n=(\epsilon_n/5)^{K_n}/(32M_n^4).
\] The number epsilon\_n/5 is below one. Outside the epsilon\_n disks about the listed zeros, C11--C13 imply \texttt{\textbar{}g(1/2+iy)\textbar{}\textgreater{}=c\_n}. Each disk cuts an interval of length at most 2epsilon\_n from the real y line. Across all m\_n cells the total excluded length is at most \[
2m_nK_n\epsilon_n=1/(32B_n).
\tag{C14}
\] Counting a disk or an overlapping interval more than once can only increase this upper bound. By C5 and C7, the retained part of {[}-R,R{]} carries Gamma polynomial mass at least \texttt{1-1/16-1/32=29/32} of the full mass, and therefore at least one half. This establishes P7 with precisely the stated constants.

\subsection{5. Lower and upper arithmetic comparisons}\label{lower-and-upper-arithmetic-comparisons}

The defining identity \texttt{2xi(s)=s(s-1)\ pi\^{}\{-s/2\}\ Gamma(s/2)\ zeta(s)} gives on the physical line \[
\frac{w_h(y)}{\sigma(y)}=
 \frac{|s|^2}{\sqrt\pi}\left|\frac{g(s)}{h(s)}\right|^2,
 \qquad s=1/2+iy.
\tag{C15}
\] This formula uses the entire removable value of g/h at selected roots. Away from h=0 it is exactly P8 in its displayed fraction form.

Let \texttt{R\_h=2+max\_rho\textbar{}rho-1/2\textbar{}}. On \textbar y\textbar\textless=R, \[
|h(1/2+iy)|\le(R+R_h)^d,\qquad |s|^2\ge1/4.
\] The retained set contains no zero of g. In particular h is nonzero there, so division is legitimate. C15 and the retained half-mass yield \[
\|P\|_{w_h}^2\ge
\frac{c_n^2}{8\sqrt\pi(R+R_h)^{2d}}\|P\|_\sigma^2
=a_{h,n}\|P\|_\sigma^2.
\tag{C16}
\]

For the upper comparison, g/h is entire because h divides the actual nontrivial-zero divisor of g. Consider the rectangle \texttt{-1\textless{}=Re\ s\textless{}=2}, \texttt{\textbar{}Im\ s\textbar{}\textless{}=2R\_h}. On its vertical sides, every \texttt{\textbar{}s-rho\textbar{}\textgreater{}=1}, since \texttt{0\textless{}Re\ rho\textless{}1}. On its horizontal sides, \[
|s-\rho|\ge2R_h-|\operatorname{Im}\rho|
 \ge2R_h-(R_h-2)=R_h+2>1.
\] Thus \textbar h\textbar\textgreater=1 on the entire boundary. Also \texttt{\textbar{}s\textbar{}+4\textless{}=2R\_h+6} there. C9 and the maximum principle applied to the entire g/h prove \[
|g/h|\le(2R_h+6)^5
\] throughout the rectangle, including all removable values. On its critical line segment, C15 consequently gives \[
w_h/\sigma\le[(2R_h)^2+1](2R_h+6)^{10}.
\tag{C17}
\]

For \textbar y\textbar\textgreater=2R\_h\textgreater=4, triangle inequalities give \[
|s-\rho|\ge|y|-|\rho-1/2|\ge|y|/2,\quad
|h(s)|\ge(|y|/2)^d\ge1,
\] and \texttt{\textbar{}s\textbar{}\textless{}=2\textbar{}y\textbar{}}, \texttt{\textbar{}s\textbar{}+4\textless{}=3\textbar{}y\textbar{}}. Hence C9 and C15 give \[
w_h/\sigma\le(4\cdot3^{10}/\sqrt\pi)|y|^{12}
 \le2^{24}(1+y^2)^6.
\tag{C18}
\] Combining C17 and C18 proves P9 with exactly \[
C_h=\max\{2^{24},[(2R_h)^2+1](2R_h+6)^{10}\}.
\] For every polynomial of degree at most m, C4 implies \texttt{\textbar{}\textbar{}(y+i)Q\textbar{}\textbar{}\_sigma\ \textless{}=(2m+3)\textbar{}\textbar{}Q\textbar{}\textbar{}\_sigma}. Iteration at \texttt{m=n,n+1,...,n+5} gives \[
\|(y+i)^6Q\|_\sigma^2\le
 \prod_{r=0}^5(2n+2r+3)^2\|Q\|_\sigma^2
 \le(2n+15)^{12}\|Q\|_\sigma^2.
\] Because \texttt{\textbar{}(y+i)\^{}6\textbar{}\^{}2=(1+y\^{}2)\^{}6}, P9 proves the upper half of P2 with the unchanged \texttt{b\_\{h,n\}=C\_h(2n+15)\^{}\{12\}}.

Now \texttt{K\_n=O(log(n+2))}, \texttt{log\ M\_n=O(log(n+2))}, and \[
-\log\epsilon_n=
\log64+\log B_n+\log m_n+\log K_n=O(\log(n+2)).
\] The explicit logarithm of a\_\{h,n\} is \[
-\log a_{h,n}=
2K_n\log(5/\epsilon_n)+2\log32+8\log M_n
+\log(8\sqrt\pi)+2d\log(R+R_h).
\tag{C19}
\] It is \texttt{O\_h(log\^{}2(n+2))}. Also \texttt{log\ b\_\{h,n\}=log\ C\_h+12\ log(2n+15)=O\_h(log(n+2))}. This proves P3 with no change to the constants.

\subsection{6. The complete divided-jet bound}\label{the-complete-divided-jet-bound}

Let \texttt{r=n/(n+1)} and \texttt{u\_n=(1/2)log(2n+1)}. On \textbar z\textbar=r the power series for arctan gives \[
|\arctan z|\le\sum_{l\ge0}\frac{r^{2l+1}}{2l+1}
=\operatorname{arctanh}r=u_n.
\] Further, \[
|1+z^2|^{-1/4}\le(1-r^2)^{-1/4}\le(n+1)^{1/4},
\qquad r^{-m}\le(1+1/n)^n\le e\quad(m\le n).
\] Write \texttt{Z\_h=max\_rho\textbar{}rho-1/2\textbar{}} and \texttt{y\_rho=(rho-1/2)/i}, so \texttt{\textbar{}y\_rho\textbar{}\textless{}=Z\_h}. Cauchy\textquotesingle s coefficient estimate applied to the j-th divided y derivative of C1 gives \[
\left|\frac{p_m^{(j)}(y_\rho)}{j!}\right|
\le m!\,e(n+1)^{1/4}(2n+1)^{Z_h/2}\frac{u_n^j}{j!}.
\tag{C20}
\] The s derivative of \texttt{p\_m((s-1/2)/i)} contributes the exact phase \texttt{i\^{}\{-j\}}, of modulus one. The inequality \texttt{m!/(1/2)\_m\textless{}=2m+1} holds by induction: its ratio at the next step is \texttt{(m+1)/(m+1/2)}, and \texttt{(2m+1)(m+1)/(m+1/2)=2m+2\textless{}=2m+3}.

Expand Q in the orthonormal basis e\_0,...,e\_n and apply Cauchy--Schwarz to its j-th divided derivative at y\_rho. C2 and C20 give \[
\left|\frac{P^{(j)}(\rho)}{j!}\right|^2
\le\|P\|_\sigma^2\frac{e^2}{M_\sigma}(n+1)^{1/2}
(2n+1)^{Z_h}\frac{u_n^{2j}}{(j!)^2}
 \sum_{m=0}^n(2m+1).
\] The sum is \texttt{(n+1)\^{}2\ \textless{}=(n+1)(2n+1)}. Thus \[
\left|\frac{P^{(j)}(\rho)}{j!}\right|^2
\le\|P\|_\sigma^2\frac{e^2}{M_\sigma}(n+1)^{3/2}
(2n+1)^{Z_h+1}\frac{u_n^{2j}}{(j!)^2}.
\tag{C21}
\] Multiply by the original weights \texttt{a\_rho\ (j!)\^{}2\ binom(d\_rho,j)} and sum over all roots and all jet orders. The factorials cancel exactly; the remaining sum over j is \texttt{(1+u\_n\^{}2)\^{}\{d\_rho\}}. Therefore \[
\|J_hP\|_{\mathsf D_{h,1}}^2\le J_{h,n}\|P\|_\sigma^2,
\quad
J_{h,n}=\frac{e^2}{M_\sigma}(n+1)^{3/2}(2n+1)^{Z_h+1}
 \sum_\rho a_\rho(1+u_n^2)^{d_\rho}.
\tag{C22}
\] Combining C22 and C16 gives P11 with \texttt{Lambda\_\{h,n\}=max\{1,J\_\{h,n\}/a\_\{h,n\}\}}. The fixed finite root set and fixed weights give \texttt{log\^{}+\ J\_\{h,n\}=O\_h(log(n+2))}; C19 therefore proves \texttt{log\ Lambda\_\{h,n\}=O\_h(log\^{}2(n+2))}.

\subsection{7. Tensor forms and exact collision coefficients}\label{tensor-forms-and-exact-collision-coefficients}

Let V\_n be the full polynomial space of degree at most n with its w\_h inner product, and let J:V\_n-\textgreater A\_h be the remainder map. C22 and C16 say that the finite operator \[
J:(V_n,\|\cdot\|_{w_h})\longrightarrow(A_h,T_1)
\] has squared operator norm at most Lambda\_\{h,n\}. Tensoring its singular value decomposition shows that the squared norm of \(J^{\otimes k}\) is at most Lambda\_\{h,n\}\^{}k. This acts on the full Hilbert tensor product \(V_n^{\otimes k}\), which is precisely the space of polynomials of separate degree at most n. Its norm is the original product-measure norm. This proves J1 for every coefficient vector, including all cross terms.

For clarity the cyclic source and its complete collision coefficients can be obtained directly. For an ordered tuple \texttt{boldrho=(rho\_1,...,rho\_k)}, put \[
\kappa(\boldsymbol\rho)=\sum_i\rho_i,\quad
D(\boldsymbol\rho)=\sum_i d_{\rho_i},\quad
A(\boldsymbol\rho)=\prod_i a_{\rho_i}.
\] Let K be the set of distinct attained kappa, and set \[
e_\kappa=1+\max_{\kappa(\boldsymbol\rho)=\kappa}D(\boldsymbol\rho),
\quad \chi(S)=\prod_{\kappa\in K}(S-\kappa)^{e_\kappa},
\quad E=\mathbb C[S]/(\chi),\quad q=\deg\chi.
\] The map \[
\beta:E\longrightarrow A_h^{\otimes k},\qquad
[P(S)]\longmapsto[P(s_1+\cdots+s_k)]
\tag{C23}
\] is well-defined and injective. In the primary factor for a tuple, \(s_i=\rho_i+t_i\) and \(t_i^{d_{\rho_i}+1}=0\). The element \texttt{t\_1+...+t\_k} has nilpotence order D+1: its D-th power is the nonzero multiple \(\displaystyle\frac{D!}{\prod_i d_{\rho_i}!}\prod_i t_i^{d_{\rho_i}}\), and every monomial of degree D+1 vanishes. For each center the largest such D therefore imposes exactly the exponent e\_kappa. Taking all tuple factors shows that the kernel of the substitution from C{[}S{]} is exactly (chi).

Write \texttt{v\_\{kappa,j\}=P\^{}\{(j)\}(kappa)/j!}. In the tuple component of beta v, the coefficient of \texttt{prod\_i\ t\_i\^{}\{j\_i\}}, with sum j\_i=j, is \[
\frac{j!}{\prod_i j_i!}v_{\kappa,j}.
\] Its squared coefficient times its exact product metric weight is \[
(j!)^2 |v_{\kappa,j}|^2 A(\boldsymbol\rho)
 \prod_i\binom{d_{\rho_i}}{j_i}.
\] The sum over the j\_i is the coefficient of x\^{}j in \texttt{prod\_i(1+x)\^{}\{d\_\{rho\_i\}\}=(1+x)\^{}D}. Thus, with zero binomial coefficients when j\textgreater D, \[
S_{\kappa,j}=\sum_{\kappa(\boldsymbol\rho)=\kappa}
A(\boldsymbol\rho)\binom{D(\boldsymbol\rho)}j,
\quad
G_{k,1}(v)=\sum_{\kappa,j}(j!)^2S_{\kappa,j}|v_{\kappa,j}|^2.
\tag{C24}
\] Each S\_\{kappa,j\} is positive on the stated j range, since a tuple of maximum D contributes positively. Every ordered tuple, including tuples with unequal lengths at a common center, is included. With D\_\{h,eta\} in every factor, the product contributes eta\^{}\{2sum j\_i\}=eta\^{}\{2j\}; hence \[
G_{k,\eta}(v)=\sum_{\kappa,j}\eta^{2j}(j!)^2S_{\kappa,j}|v_{\kappa,j}|^2
\preceq G_{k,1}(v).
\tag{C25}
\] For every full polynomial F with \texttt{J\^{}\{tensor\ k\}F=beta\ v}, J1 and C24 give the exact J2 lower bound \texttt{\textbar{}\textbar{}F\textbar{}\textbar{}\^{}2\textgreater{}=Lambda\_\{h,n\}\^{}\{-k\}G\_\{k,1\}(v)}.

\subsection{8. The section and the complete minimum}\label{the-section-and-the-complete-minimum}

Set n\_0=2d-1, and use the original coefficient frame on V\_\{n\_0\}. Its source Gram H is positive definite: w\_h is positive except at a discrete set, so the squared norm of a nonzero polynomial is positive. The remainder matrix J has rank d. Its kernel consists of \texttt{hQ} with deg Q \textless= d-1, so it has dimension d. Define \[
G_0=(JH^{-1}J^*)^{-1},\quad R_0=H^{-1}J^*G_0.
\] Then \texttt{JR\_0=I} and \texttt{R\_0\^{}*HR\_0=G\_0}. Choose W with d columns spanning ker J and \texttt{W\^{}*HW=I\_d}. Direct multiplication gives \texttt{R\_0\^{}*HW=G\_0JW=0}.

Let \texttt{L\_h=1+\textbar{}\textbar{}T\_1\^{}\{-1/2\}G\_0T\_1\^{}\{-1/2\}\textbar{}\textbar{}} and \texttt{tau\_eta=L\_h\ eta\^{}\{-2d\_max\}}. The exact jet weights imply \[
\eta^{2d_{\max}}T_1\preceq T_\eta\preceq T_1,
\qquad \tau_\eta T_\eta\succeq L_hT_1\succ G_0.
\] The positive square root in the original coefficient frame is consequently well-defined. With \texttt{D=(tau\_eta\ T\_eta-G\_0)\^{}\{1/2\}}, put \[
R_\eta=R_0+WD.
\] The exact equalities are \[
JR_\eta=I,\qquad R_\eta^*HR_\eta
=G_0+D^*D=\tau_\eta T_\eta.
\tag{C26}
\] No square root has been transported through a nonunitary congruence.

Tensor C26 and compose with beta. After the original physical source isometry A\_h, the section \texttt{mathcal\ R=(A\_h\ R\_eta)\^{}\{tensor\ k\}\ beta} has value beta v and squared norm \[
\mathcal R^*\mathcal R=\tau_\eta^kG_{k,\eta}.
\tag{C27}
\] Let X\_N be exactly the scalar-polynomial physical source from P(S), deg P\textless=N, with S=s\_1+...+s\_k, and set Y=im mathcal R. Both sources lie in the box of separate degree at most \texttt{n=max(N,n\_0)}. The section shows that their sum maps onto E. Its minimum exists and is unique modulo zero-value directions by finite-dimensional orthogonal projection. Its metric is positive definite. J2 applied before taking this minimum gives its lower bound; C27 is an admitted vector for each value and gives its upper bound. Thus \[
\Lambda_{h,n}^{-k}G_{k,1}\preceq G_N^J
\preceq\tau_\eta^kG_{k,\eta}\preceq\tau_\eta^kG_{k,1}.
\tag{C28}
\] For N\textgreater=q-1 the scalar source is onto E, with positive minimum metric G\_N; source inclusion then gives \texttt{G\_N\^{}J\textless{}=G\_N}.

For the full box, let H\_n be the one-factor polynomial Gram and J\_n its remainder matrix. Its covariance is \texttt{C\_n=J\_n\ H\_n\^{}\{-1\}\ J\_n\^{}*}, and its minimum metric is \texttt{G\_\{h,n\}=C\_n\^{}\{-1\}}. The box covariance is exactly \[
(J_n^{\otimes k})(H_n^{-1})^{\otimes k}(J_n^*)^{\otimes k}
=C_n^{\otimes k}.
\] Its value metric is therefore \texttt{G\_\{h,n\}\^{}\{tensor\ k\}}. Restriction to the exact value beta v yields \[
G_{k,n}^{\rm box}=\beta^*(G_{h,n}^{\otimes k})\beta\preceq G_N^J,
\tag{C29}
\] because every joint candidate is a box candidate. This proves J4 with the whole box kernel present. The same proof factor by factor gives the heterogeneous constants \texttt{prod\_i\ Lambda\_\{h\_i,n\_i\}} and \texttt{prod\_i\ tau\_\{i,eta\}}, with \texttt{n\_i=max(N,2d\_i-1)}.

\subsection{9. The original cross Gram and its positive-range inverse}\label{the-original-cross-gram-and-its-positive-range-inverse}

For N\textgreater=q-1 let B\_X be the scalar coefficient map into the actual source, \texttt{H\_X=B\_X\^{}*B\_X}, and let J\_N be its map onto E. Its orthogonal projector is \texttt{P\_X=B\_X\ H\_X\^{}\{-1\}B\_X\^{}*}. Put \[
C_X=B_X^*\mathcal R,\quad
R_Z=(I-P_X)\mathcal R,\quad
H_Z=R_Z^*R_Z=\Gamma-C_X^*H_X^{-1}C_X,
\] \[
F_Z=I-J_NH_X^{-1}C_X.
\tag{C30}
\] The actual value map sends R\_Z to beta F\_Z. R\_Z lies in the same polynomial box, so J1 applied to R\_Z v for every v gives \[
F_Z^*G_{k,1}F_Z\preceq\Lambda_{h,n}^kH_Z.
\tag{C31}
\] In particular \texttt{ker\ H\_Z\ subset\ ker\ F\_Z}, since G\_\{k,1\} is positive.

The sum source is the orthogonal sum \texttt{X\_N\ +\ im\ R\_Z}. Take a unitary spectral decomposition of H\_Z. On each positive eigenvalue lambda its source coordinate becomes orthonormal after multiplication by lambda\^{}\{-1/2\}; coordinates in ker H\_Z represent the zero vector and are killed by F\_Z. Hence the covariance of this normal source is exactly \texttt{F\_Z\ H\_Z\^{}+\ F\_Z\^{}*}. The scalar covariance is \texttt{J\_N\ H\_X\^{}\{-1\}\ J\_N\^{}*=G\_N\^{}\{-1\}}. Covariances of orthogonal source summands add, which proves \[
\Omega_N=F_ZH_Z^+F_Z^*,\qquad
G_N^J=(G_N^{-1}+\Omega_N)^{-1}.
\tag{C32}
\] The inverse exists because the scalar covariance is positive definite. Factoring it out gives \[
\frac{\det G_N}{\det G_N^J}
=\det(I+G_N^{1/2}\Omega_NG_N^{1/2}),
\tag{C33}
\] which proves J7--J8 on the original coefficient space.

The section is injective by its right-inverse property. Therefore \[
\dim\ker R_Z=\dim(Y\cap X_N),\quad
\operatorname{rank}H_Z=q-\dim(Y\cap X_N).
\] If a vector belongs to the intersection, its polynomial is both P(s\_1+...+s\_k) and of separate degree at most n\_0. Equality of physical vectors implies equality of these polynomials, because the product density is positive almost everywhere and a polynomial vanishing there is zero. If deg P=M, its highest pure s\_1 monomial has nonzero coefficient and degree M, so M\textless=n\_0. Consequently \[
\dim(Y\cap X_N)\le\min(N,n_0)+1\le2d,
\qquad \operatorname{rank}H_Z\ge q-2d.
\tag{C34}
\] This is J9. It makes no claim that the value correction Omega\_N has the same rank as H\_Z.

\subsection{10. Fixed restrictions, complete quotients, and the four cutoffs}\label{fixed-restrictions-complete-quotients-and-the-four-cutoffs}

Assume \texttt{2q\textgreater{}=n\_0} and use the single constant \texttt{Lambda=Lambda\_\{h,2q\}} throughout \texttt{q-1\textless{}=N\textless{}=2q}. Every relevant source lies in that same box. With \texttt{G\_ref=G\_\{k,1\}}, C28 gives \[
\alpha G_{\rm ref}\preceq G_N^J\preceq bG_{\rm ref},
\quad \alpha=\Lambda^{-k},\quad b=\tau_\eta^k,
\quad \log(b/\alpha)=W_k=k\log(\tau_\eta\Lambda).
\tag{C35}
\] For any fixed injective inclusion I, congruence gives the same bounds on \texttt{I\^{}*G\_N\^{}J\ I}. For any fixed onto quotient L onto its actual image, minimizing the scalar quadratic inequality over the exact affine fibre \texttt{Lv=w} gives the same bounds on its quotient metric. That minimum is \[
Q_N=(L(G_N^J)^{-1}L^*)^{-1},
\] as follows by the minimizing vector \texttt{v=(G\_N\^{}J)\^{}\{-1\}L\^{}*Q\_N\ w}, whose difference from any other lift is orthogonal to ker L. Repeated fixed inclusions and quotients preserve these same constants, proving the assertion for every fixed subquotient.

Because Y is fixed and X\_N increases, G\_N\^{}J decreases with N. Restriction and affine minimization preserve this order. If the resulting rank is r\_X and \texttt{f(N)=log\ det\ G\_\{X,N\}\^{}J}, positive congruence and eigenvalue comparison imply \[
f(N_1)\ge f(N_2)\quad(N_1\le N_2),\qquad
\sup_Nf(N)-\inf_Nf(N)\le r_XW_k.
\] Pair the actual four signs as \[
\mathcal Rf=[f(q-1)-f(2q-1)]+[f(q)-f(2q)].
\] Both differences lie between zero and r\_X W\_k. This proves J11 exactly: \[
0\le\mathcal R\log\det G_{X,N}^J\le2r_XW_k.
\tag{C36}
\] Rank one includes the restriction to a specified nonzero class, with its original vector fixed across all cutoffs.

The complete finite tolerance expression is \[
W_k=k\log L_h+k\log\Lambda+2kd_{\max}\log(1/\eta).
\tag{C37}
\] For the general choice \texttt{eta=min(1,t\_*/(D+1))}, where D is the largest total tuple degree and \texttt{D\textless{}=k\ d\_max}, \[
\log(1/\eta)=\log^+((D+1)/t_*)
 \le\log(D+1)+\log^+(1/t_*).
\] C19 and C22 therefore give the precise J10 order \[
W_k=O_h\bigl(k\log^2(q+2)+k\log(k+2)
 +k\log^+(1/t_*)\bigr).
\tag{C38}
\] For a common-multiplicity quartet with \texttt{e=1+k(m\_0-1)} and \texttt{eta=delta/e} when m\_0\textgreater1, one instead reads \texttt{log(1/eta)=log\ e+log(1/delta)=O\_h(log(k+2))} directly. For simple roots d\_max=0, so C37 is independent of eta and gives \texttt{W\_k=O\_h(k\ log\^{}2(q+2))}; on the simple quartet \texttt{q=(k+1)\^{}2}.

For every fixed one-factor root set with z=\textbar Z\textbar, there is also the explicit polynomial growth bound \[
q\le(kd_{\max}+1)\binom{k+z-1}{z-1}.
\tag{C39}
\] Indeed an unordered root-count vector is a z-tuple of nonnegative integers summing to k, so there are exactly the displayed binomial number of such vectors. Every ordered tuple has one such vector and therefore one of at most that many centers. Each center has e\_kappa\textless=k d\_max+1. Summing these inequalities proves C39, including all possible collisions. This is the polynomial growth assertion needed in J10. On the stated quartet the exact formula q=(1+k(m\_0-1))(k+1)\^{}2 follows because the independent real and imaginary sign counts each run from 0 through k and every pair of counts is attainable; all tuples have the same D=k(m\_0-1).

If a family of intermediate sources contains Y and stays inside the same box, C35 and its fixed restrictions and quotients still hold. Nesting is needed for the first inequality in C36. In the absence of nesting each of the two paired differences lies in \texttt{{[}-r\_XW\_k,r\_XW\_k{]}}, yielding the valid conclusion \texttt{\textbar{}R\ f\textbar{}\textless{}=2r\_X\ W\_k}. Common enclosing constants alone cannot fix the sign: in rank one, choosing both early metric values at alpha and both late values at b gives \texttt{R\ f=-2log(b/alpha)}. This shows exactly why the claimed intermediate-source extension needs its stated scope repair. The original joint source satisfies the nesting requirement.

\subsection{11. Reading and use record}\label{reading-and-use-record}

The mathematical input actually read for this derivation was:

\begin{itemize}
\tightlist
\item
  \path{COMPLETE_PROOFS(2).md}, complete Sections 1--2, equations P1--P11 and J1--J11. The initial tool output also exposed later sections; no independent verification of those later sections is asserted.
\item
  \path{.../family_cost_intake_20260920/extracted/ES_RH_Family_Integration_20260920_C/proofs/01_METRIC_RETURN.md}, complete file read; only Section 1.1\textquotesingle s explicit common-multiplicity tolerance choice was used here. Its separate imported spectral theorem was not used in the present proof.
\item
  The same package\textquotesingle s \path{proofs/02_JOINT_SOURCE.md}, bounded search excerpts around the cross-Gram and physical-source discussion. The cross-Gram identities in this file were independently proved in C30--C33; no claim of complete reading of that package is made.
\item
  NIST DLMF, version 1.2.8, release 2026-09-15: exact formulas \href{https://dlmf.nist.gov/5.7.E6}{5.7.6}, \href{https://dlmf.nist.gov/18.23.E7}{18.23.7}, \href{https://dlmf.nist.gov/18.19.E8}{18.19.8--9}, and \href{https://dlmf.nist.gov/25.2.E10}{25.2.10--11}. These formulas were read in the rendered primary source and fetched from its original \path{.tex} endpoints. They supply only the classical identities identified in C1--C3, C6, and C8; all quantitative constants and all finite tensor and minimum arguments are derived above. The six unchanged TeX responses are retained in \texttt{coercivity\_sources/} under their DLMF equation IDs, including \path{18.19.E9a.tex} for the norm.
\end{itemize}

No paper absent from this list, no original arXiv source, and no entire literature corpus is represented as having been read. No PDF was used. No remote publication or graphical rendering was performed in this independent derivation subtask.

\end{document}
